% Slide build of Figure 21.1 -- resolution one hop at a time.
%
% The whole point of the figure is that step 1 is recursive and steps 2-4 are
% iterative: the host asks once and waits, the resolver does the walking. Drawn
% all at once, that distinction is a thing the lecturer asserts. Built up, the
% room watches the host go quiet after step 1 while the resolver runs around,
% and the assertion is unnecessary.
%
% Step 5 adds the cache, because "and it remembers for the TTL" is the answer
% to "why is it fast the second time" and to half of chapter 21's Break & Fix.
\begin{tikzpicture}[font=\scriptsize]
  % Fixed canvas: the widest step must not resize the picture, or Morph
  % translates the whole diagram instead of fading the new arrow in.
  \useasboundingbox (-2.1,-4.3) rectangle (11.9,2.7);

  \tikzset{
    bx/.style={rounded corners=3pt, draw=#1, fill=#1!8, text=#1!45!black,
               line width=0.9pt, text width=2.5cm, align=center,
               minimum height=1.15cm, font=\scriptsize\bfseries},
    q/.style={-{Stealth[length=2.2mm]}, line width=0.9pt, dom4!75},
    a/.style={-{Stealth[length=2.2mm]}, line width=0.9pt, dom1!70,
              dash pattern=on 2.5pt off 2pt},
    stepcap/.style={font=\bfseries, anchor=north west, text width=12.6cm,
                    align=left}}

  \node[bx=neutral] (h)  at (0,0)     {\faIcon{desktop}\\[1pt]Host\\
      \footnotesize\mdseries stub resolver};
  \node[bx=dom1]    (rec) at (4.9,0)  {\faIcon{server}\\[1pt]Recursive\\
      \footnotesize\mdseries resolver, caches};
  \node[bx=dom4]    (rt)  at (9.7,1.5){\faIcon{globe}\\[1pt]Root\\
      \footnotesize\mdseries ``ask .edu''};
  \node[bx=dom4]    (tld) at (9.7,0)  {\faIcon{globe}\\[1pt]TLD .edu\\
      \footnotesize\mdseries ``ask them''};
  \node[bx=dom4]    (au)  at (9.7,-1.5){\faIcon{database}\\[1pt]Authoritative\\
      \footnotesize\mdseries holds the zone};

  % Query below the line, answer above the arc that bends over it. With both
  % on the inside they printed on top of each other in the gap between the
  % nodes -- which the book did too, and is now fixed in both.
  \stepvis{1}{\draw[q] (h) -- node[font=\tiny, below, midway,
      text=dom4!45!black] {1. recursive query} (rec);}
  \stepvis{2}{\draw[q] (rec) -- node[font=\tiny, above=1pt, pos=0.6] {2} (rt);}
  \stepvis{3}{\draw[q] (rec) -- node[font=\tiny, above=1pt, pos=0.6] {3} (tld);}
  \stepvis{4}{\draw[q] (rec) -- node[font=\tiny, above=1pt, pos=0.6] {4} (au);}
  \stepvis{5}{\draw[a] (rec) to[bend right=22]
      node[font=\tiny, above, midway, text=dom1!45!black]
      {5. answer, with a TTL} (h);}

  % The host is doing nothing at all during 2-4. Saying so on the diagram
  % is worth more than saying it aloud -- and it has to come off again at
  % step 5, when the host is no longer waiting for anything.
  \stepvis{2}{\stepdim{5}{\node[font=\tiny\itshape, text=neutral!70!black,
      anchor=north] at (0,-1.0) {still waiting};}}

  \stepvis{5}{%
    \node[draw=dom1!55, fill=dom1!8, rounded corners=3pt, line width=0.8pt,
          font=\tiny\ttfamily, align=left, inner sep=5pt, anchor=north west]
        at (3.4,-1.9)
        {\textcolor{dom1!45!black}{cache}\\ www.example.edu\\ TTL 3600};}

  \steponly{1}{\node[stepcap, text=dom4!45!black] at (-2.0,-3.0)
      {One question, then the host waits. This arrow is the only recursive
       one in the picture.};}
  \steponly{2}{\node[stepcap, text=dom4!45!black] at (-2.0,-3.0)
      {The root does not know the answer --- it knows who runs .edu.};}
  \steponly{3}{\node[stepcap, text=dom4!45!black] at (-2.0,-3.0)
      {Nor does the TLD. Each referral is one step down the name, right
       to left.};}
  \steponly{4}{\node[stepcap, text=dom4!45!black] at (-2.0,-3.0)
      {The authoritative server holds the zone, so this one answers.};}
  \steponly{5}{\node[stepcap, text=dom1!45!black] at (-2.0,-3.0)
      {The host gets one answer and never learns any of that happened ---
       and the resolver keeps it for the TTL, so the next host skips
       steps 2 to 4.};}
\end{tikzpicture}
